% Options for packages loaded elsewhere
\PassOptionsToPackage{unicode}{hyperref}
\PassOptionsToPackage{hyphens}{url}
\documentclass[
  a4paper,
]{ltjsarticle}
\usepackage{xcolor}
\usepackage[margin=25mm]{geometry}
\usepackage{amsmath,amssymb}
\setcounter{secnumdepth}{-\maxdimen} % remove section numbering
\usepackage{iftex}
\ifPDFTeX
  \usepackage[T1]{fontenc}
  \usepackage[utf8]{inputenc}
  \usepackage{textcomp} % provide euro and other symbols
\else % if luatex or xetex
  \usepackage{unicode-math} % this also loads fontspec
  \defaultfontfeatures{Scale=MatchLowercase}
  \defaultfontfeatures[\rmfamily]{Ligatures=TeX,Scale=1}
\fi
\usepackage{lmodern}
\ifPDFTeX\else
  % xetex/luatex font selection
\fi
% Use upquote if available, for straight quotes in verbatim environments
\IfFileExists{upquote.sty}{\usepackage{upquote}}{}
\IfFileExists{microtype.sty}{% use microtype if available
  \usepackage[]{microtype}
  \UseMicrotypeSet[protrusion]{basicmath} % disable protrusion for tt fonts
}{}
\makeatletter
\@ifundefined{KOMAClassName}{% if non-KOMA class
  \IfFileExists{parskip.sty}{%
    \usepackage{parskip}
  }{% else
    \setlength{\parindent}{0pt}
    \setlength{\parskip}{6pt plus 2pt minus 1pt}}
}{% if KOMA class
  \KOMAoptions{parskip=half}}
\makeatother
\usepackage{longtable,booktabs,array}
\newcounter{none} % for unnumbered tables
\usepackage{calc} % for calculating minipage widths
% Correct order of tables after \paragraph or \subparagraph
\usepackage{etoolbox}
\makeatletter
\patchcmd\longtable{\par}{\if@noskipsec\mbox{}\fi\par}{}{}
\makeatother
% Allow footnotes in longtable head/foot
\IfFileExists{footnotehyper.sty}{\usepackage{footnotehyper}}{\usepackage{footnote}}
\makesavenoteenv{longtable}
\usepackage{graphicx}
\makeatletter
\newsavebox\pandoc@box
\newcommand*\pandocbounded[1]{% scales image to fit in text height/width
  \sbox\pandoc@box{#1}%
  \Gscale@div\@tempa{\textheight}{\dimexpr\ht\pandoc@box+\dp\pandoc@box\relax}%
  \Gscale@div\@tempb{\linewidth}{\wd\pandoc@box}%
  \ifdim\@tempb\p@<\@tempa\p@\let\@tempa\@tempb\fi% select the smaller of both
  \ifdim\@tempa\p@<\p@\scalebox{\@tempa}{\usebox\pandoc@box}%
  \else\usebox{\pandoc@box}%
  \fi%
}
% Set default figure placement to htbp
\def\fps@figure{htbp}
\makeatother
\setlength{\emergencystretch}{3em} % prevent overfull lines
\providecommand{\tightlist}{%
  \setlength{\itemsep}{0pt}\setlength{\parskip}{0pt}}
\usepackage{bookmark}
\IfFileExists{xurl.sty}{\usepackage{xurl}}{} % add URL line breaks if available
\urlstyle{same}
\hypersetup{
  hidelinks,
  pdfcreator={LaTeX via pandoc}}

\author{}
\date{}

\begin{document}

\section{空間軸3軸と固有時間の創生------凝縮体からの読出しの再導出と、唯一読めない座標時間}\label{ux7a7aux9593ux8ef83ux8ef8ux3068ux56faux6709ux6642ux9593ux306eux5275ux751fux51ddux7e2eux4f53ux304bux3089ux306eux8aadux51faux3057ux306eux518dux5c0eux51faux3068ux552fux4e00ux8aadux3081ux306aux3044ux5ea7ux6a19ux6642ux9593}

\textbf{著者:} 木原 範昭（WF System Co., Ltd.）　\textbf{日付:}
2026-08-06　\textbf{版:} v1 \textbf{Version DOI:}
\href{https://doi.org/10.5281/zenodo.21816652}{10.5281/zenodo.21816652}
\textbf{Concept DOI:}
\href{https://doi.org/10.5281/zenodo.21816651}{10.5281/zenodo.21816651}

\begin{center}\rule{0.5\linewidth}{0.5pt}\end{center}

\subsection{要旨}\label{ux8981ux65e8}

先行論文は、物質パケットの形不変並進を実証したが、その位置軸（3次元レジスタ格子）と時計（調和閉鎖分散）は外から与えていた。本稿はその両方を、三方向準安定凝縮体自身の観測量から\textbf{導出し直す}。あわせて質量的な量と電荷的な量の読出しも同じ枠で再導出し、読出し方法をすべて本文で明記する。実測の結果：(1)
凝縮体の実体は対等な3軸ではなく \textbf{2＋1}
である------独立な実体は複素固有対1つ（回転平面1枚）であり、第3方向は生成子の回転軸として従属する（外積との一致
0.987→0.9965、Nとともに向上）。(2) 空間の \(x, y\)
軸は、回転平面の複素座標を固有時計で復調した\textbf{直交位相成分（quadrature）}として導出される------直交性は実体の性質ではなく読出しの構成である。(3)
固有時間 τ
は凝縮体の\textbf{単一の集団時計}（ω≈π/72/step、全占有平面が剛体回転）の刻みであり、この時計は\textbf{引き込み}によって集団強制される：時計を持たないコヒーレント摂動は凝縮体の時計に同期させられて静止する。静止成分への変位は無減衰で保持され、運動は自前の時計を持つ内容------質量------の特権となる。(4)
質量は波形の\textbf{非コヒーレンス}
\(\det\Gamma = T^2\!-\!X^2\!-\!Y^2\!-\!Z^2 \ge 0\)
として読める（数学は偏光のコヒーレンス行列理論として既知であり、本稿の寄与は物理的同定と模型内実測である：光的固有モード
\(1.2\times10^{-7}\) 対 凝縮体 \(3.1\times10^{-3}\)）。(5)
電荷は透過率・相対位相・巻き数という波形構造の読出しで分離できる。(6)
導出した読出しの\textbf{分布}を実測した：海のモードは成長期・熱平衡期とも最大非コヒーレンス（質量²/T²≈1）に集中する------海は「重い」。電荷（巻き数）は複数の量子化値に分布して±1に安定せず------既知の未解決課題として明記する。対称な海のフェルミオン的内容は灰色（粒子・反粒子均衡）に、帯電した種を持つ構造は粒子的に集中する。物質パケットの運動は
xyz＋τ
の4次元読出しで等速並進として可視化される（τ=0〜900の全10点で偏差0.000）。以上を閉塞の錐
\(x^2+y^2+z^2 = t^2+R^2+Q^2 = R'^2\)
の上に整理すると、\textbf{観測の全内容は空間3軸＋固有時間τ＋質量m＋電荷qであり、唯一読めないのは座標時間
t である}------t は残差 \(t^2 = R'^2 - m^2 - q^2\)
として、または大きな系の空間振動群から広域の共有時間を想定する選択としてのみ現れる。帰結として、一般相対論が公理的に置く「固有時間は物理・座標時間は規約」の区別が、読出し構造から導出される。次元の鋭さは有限で（曖昧さ
\(\sim 1/M\)）、N=4では第3次元が結晶化しない------少数体複合粒子（バリオン）の領域は本稿の対象外として切り出す。

\begin{center}\rule{0.5\linewidth}{0.5pt}\end{center}

\subsection{1.
本稿の性格------前提・導出・実測の三層分離}\label{ux672cux7a3fux306eux6027ux683cux524dux63d0ux5c0eux51faux5b9fux6e2cux306eux4e09ux5c64ux5206ux96e2}

本稿は、新しい力学を提案する論文ではない。三層を分離して述べる。

\begin{itemize}
\tightlist
\item
  \textbf{(a) 前提}：継承する力学（第7・8論文の線形力学、read-only
  import・SHA記録）と、三方向準安定凝縮体（第7論文の発見、先行論文の通し走行で創生過程まで実測済み）。
\item
  \textbf{(b) 導出}：\(x, y, z\)・固有時間 τ・質量 m・電荷 q
  の読出し方法の構成（§6--§9）。\textbf{引用は出所表示であって導出の代用ではない}------全読出しを本文で再導出し、本稿単体で検証が閉じるように書く。
\item
  \textbf{(c) 実測}：予備実験系列
  v1--v5・MC（実験一覧・全結果JSON・対照検定つき）。
\end{itemize}

\subsection{2. 主張}\label{ux4e3bux5f35}

\textbf{第一に------本稿の目的------読出しの再導出。} 空間3軸
\(x,y,z\)、固有時間 τ、質量的な量 m、電荷的な量 q
の読出し方法を、レジスタ格子も分散時計も手で置かずに、三方向準安定凝縮体の観測量だけから再度導出し直した。各読出しの導出は本文で明記する（§6--§9）。

\textbf{第二に、実体の構造（実測、§5）。}
凝縮体の実体は対等な3軸ではなく \textbf{2＋1}
である：占有の最上位は複素固有対1つ（回転平面1枚、接線写像の
\(|\lambda|=0.99998\)
の共役対）であり、第3方向は独立自由度ではなく、生成子を占有部分空間へ射影した反対称流の\textbf{回転軸}------上位2方向の外積との一致は
0.987（N=5）→ 0.9965（N=8）で、Nとともに向上する。

\textbf{第三に、空間3軸の読出しの導出（§6）。} \(x\) と \(y\)
は、回転平面の複素座標 \(c(n)\) を固有時計位相 \(\varphi(n)\)
で復調した直交位相成分として導出される：\(x=\mathrm{Re}\,(c\,e^{-i\varphi})\),
\(y=\mathrm{Im}\,(c\,e^{-i\varphi})\)。\(z\)
は回転軸方向の復調射影である。\textbf{直交性は実体の性質ではなく、読出しの構成である。}
先行論文で外から与えた3次元レジスタは、この内在読出しの外部代用だった。

\textbf{第四に、固有時間 τ の読出しの導出（§7）。} τ は
R′時計------凝縮体の単一の集団時計------の刻みである。実測：全占有平面が単一周波数
ω≈π/72/step（力学の基本角と0.04\%一致）で剛体回転する。さらにこの時計は\textbf{引き込み}で集団強制される：時計を持たないコヒーレント摂動は固有の離調を保持できず凝縮体の時計に同期する。静止＝時計ロックであり、静止成分への変位は無減衰で保持される。運動は自前の時計を持つ内容------質量------の特権であり、質量側の実証は先行論文の物質パケット並進が既に与えている。その全体は
xyz＋τ
の4次元読出しで等速並進として可視化される（図8：3次元軌道の等間隔点列と
x--τ, y--τ, z--τ 3投影面の直線、τ=0〜900の全10点で偏差0.000）。

\textbf{第五に、質量的な量 m の読出しの再導出（§8）。} m は R
成分の大きさであり、波形の構造から3系統で読める：(i)
非コヒーレンス------読出し対の Gram 行列の行列式
\(\det\Gamma = T^2-X^2-Y^2-Z^2\)（本文で全段導出。数学は偏光理論として既知、同定が本稿の寄与）、(ii)
閉鎖完走コスト、(iii) 自前時計の強度。

\textbf{第六に、電荷的な量 q の読出しの再導出（§9）。} q は Q
成分の大きさ＋相対符号であり、波形の構造から3系統で読める：(i)
透過率、(ii) 相対位相（電荷文法）、(iii) 巻き数（KK型運動量）。

\textbf{第七に、読出しの分布（実測の補遺、§9.5）。}
導出した読出しをモード集合 \((k,m)\)
に適用し、三つの状態（対称な海の成長期・熱平衡期・帯電した構造）で分布を実測した。(a)
海のモードはどちらの時期も\textbf{最大非コヒーレンス（質量²/T²≈1）に集中}する------海は重く、軽い（光的）状態はコヒーレントな構造に限られる。(b)
\textbf{電荷（巻き数）は複数の量子化値に分布し、±1
に安定しない}------素電荷の一意性の起源は既知の未解決課題であり、本稿はその分布の直接実測を与える（§15-6）。(c)
スピン的な量（Gram の Bloch
ベクトル------探索的読出し）では海はほぼ無偏極。(d)
フェルミオン的な波は符号均衡により\textbf{粒子的・反粒子的・灰色}に三分類でき、対称な海は灰色に、帯電した種を持つ構造は粒子的に集中する------粒子/反粒子の区別は対称な海には存在せず、帯電した種が持ち込む。

\textbf{第八に------中心命題------観測の全内容と、唯一読めないもの（§10）。}
閉塞の錐 \(x^2+y^2+z^2 = t^2+R^2+Q^2 = R'^2\)
において、観測の全内容は\textbf{空間3軸＋固有時間τ＋質量m＋電荷q}である。\textbf{唯一読めないのは座標時間
t である}：t の直接読出し器は存在せず、t は算術残差 \(t^2=R'^2-m^2-q^2\)
として、または大きな系の空間振動群から広域の共有時間を想定する\textbf{選択}としてのみ現れる。帰結は二つ。(1)
一般相対論が公理的に置く「固有時間は物理・座標時間は規約」の区別が、読出し構造から\textbf{導出}される。(2)
観測される時空がなぜ 3＋1
かへの答えの候補------読出し可能な内容がちょうど実3軸＋時計1本だから。

\textbf{第九に、向きの規約（§11）。}
時間の矢は存在しない。虚数符号は「この軸の向きは読めない」という記帳であり、閉塞は二乗しか見ない。物理量は相対符号のみであり、時計×回転軸の相対掌性（−）が規約不変量として安定に実測される。

\textbf{第十に、適用域の境界（§12）。} 第3次元の鋭さは有限である：曖昧さ
\(1-|\hat n\cdot\hat a| \sim 1/M\)。N=4
では第3次元が結晶化しない。ただし N=4,5
は少数体複合粒子（バリオン）の領域であり、次元・時間・存在が曖昧なのはその領域の物理である------本稿の対象は空間が存在する大N側であり、少数体領域は別の研究主題として切り出す。

\textbf{第十一に、反証と修正の全過程を記録する（§13）。}
計器の誤り5件------うち1件は「分解能の限界が構造をでっち上げる」実例であり、本稿の主題が計器側に現れた。

\textbf{最後に、本稿が主張しないことと反証条件（§16）。}

\subsection{3.
動機------外から与えていた二つのもの}\label{ux52d5ux6a5fux5916ux304bux3089ux4e0eux3048ux3066ux3044ux305fux4e8cux3064ux306eux3082ux306e}

先行論文（多体接続論文
{[}1{]}）は、真空→インフレーション→凝縮→物質析出をタイミング則なしの単一力学で通しで走らせ、生成された物質パケットが形を変えずに等速で動くことまでを実証した。しかしその運動デモには、手で置いたものが二つあった：\textbf{位置の軸}（3次元レジスタ格子
\(8^3\)）と、\textbf{時計}（調和閉鎖分散 \(\omega_k = k\,\omega_1\)
{[}6{]} をレジスタ時計として挿入）である。

空間と時間を外から与えたままでは、「空間軸3軸と固有時間の創生」を語ったことにならない。本稿の問いはこうである：\textbf{凝縮体は、自分の観測量だけから、位置の軸と時計を供給できるか。}
答えは肯定であり、しかも予想より強い形------3軸は対等ではなく2＋1であり、直交性は読出しが作り、時計は1本で引き込みにより集団強制され、そして読めない量がただ一つ残る------だった。

\subsection{4.
前提------力学と凝縮体}\label{ux524dux63d0ux529bux5b66ux3068ux51ddux7e2eux4f53}

力学は第7・8論文 {[}3,4{]} の原本コードを無改変で継承する（read-only
import・SHA記録済み系列）。\(N\) 体の関係波
\(Z \in \mathbb{C}^M\)（\(M=N(N-1)/2\)）、生成子
\(K_{ef}=\sin(\theta_f-\theta_e)\)（隣接辺・反対称）、Cayley一歩。自己無撞着閉包から出発した系は潜伏・急拡大ののち\textbf{三方向準安定構造}（第7論文の発見）に達する。本稿の全測定はこの準安定凝縮体の上で行う（整定後の窓、実験一覧参照）。以下、力学の\textbf{更新回数（計算上のステップラベル）を
\(n\)} と書く------これは閉塞の錐に現れる座標時間 \(t\)
とは別物である（\(t\)
は§10の主題であり、本稿では錐の分解の中にしか現れない）。親平面（自己無撞着解
\(v\) の実部・虚部が張る平面）の複素座標の偏角を \(\varphi(n)\)
と書き、親平面を複素射影で除いた成分を \(Z_\perp\) と書く。

\subsection{5.
実体の構造------2＋1（実測）}\label{ux5b9fux4f53ux306eux69cbux902021ux5b9fux6e2c}

準安定窓の軌道に対する三つの測定が、凝縮体の実体を確定する（図1）。

第一に、\(Z_\perp\) の軌道の特異値分解は\textbf{厳密に対で}並ぶ（1.000,
0.998 / 0.427, 0.426 / 0.207, 0.207 /
\ldots）------占有は回転平面の階層である。第二に、一段写像の接線行列（\(2M\times2M\)、有限差分）の固有分解で、占有部分空間を担うのは\textbf{複素共役固有対1つ}（\(|\lambda|=0.999984\)）である------独立な実体は回転平面1枚であり、第3方向に独立な固有値は存在しない。第三に、生成子
\(K\) を占有3次元部分空間へ射影した反対称 \(3\times3\) の回転軸
\(\hat a\) は、平面の2方向の外積 \(\hat n = \hat d_1\times\hat d_2\)
と一致する：\(|\hat n\cdot\hat a| = 0.9874\)（N=5）→ 0.9918（N=6）→
0.9965（N=8）。

\begin{figure}
\centering
\pandocbounded{\includegraphics[keepaspectratio,alt={図1 実体の2+1構造}]{figs/fig_s1_structure_v1.png}}
\caption{図1 実体の2+1構造}
\end{figure}

\textbf{図1}：(a) 特異値は厳密に対で並ぶ＝回転平面の階層。(b)
接線写像の固有値：占有を担うのは複素共役対1つ（赤）。第3方向は独立固有値を持たない。

まとめると、実体は\textbf{平面1枚（複素1自由度）＋従属な回転軸1本}------「三方向」の実体は
3 ではなく 2＋1、より正確には複素1＋実1である。

\subsection{6.
空間3軸の読出しの導出}\label{ux7a7aux95933ux8ef8ux306eux8aadux51faux3057ux306eux5c0eux51fa}

\textbf{導出。}
平面内を回転する状態は、平面内の\textbf{どの}2方向への射影も完全相関する（実測：読出し相関
1.0000）------ゆえに平面は独立な実読出しを2つ持たない。平面が持つ独立な読出しは複素座標1つ：

\[c(n) = \langle d_{\text{平面}},\, Z_\perp(n)\rangle\]

である（\(d_{\text{平面}}\)
は複素固有対から構成した平面の複素方向）。これを固有時計位相
\(\varphi(n)\)（§7）で復調し、\textbf{直交位相成分（quadrature）に分解}する：

\[x(n) = \mathrm{Re}\,\big(c\,e^{-i\varphi}\big), \qquad y(n) = \mathrm{Im}\,\big(c\,e^{-i\varphi}\big).\]

ここで重要な点を明示する：\(\varphi(n)\) は \(c(n)\)
自身の偏角\textbf{ではない}（もしそうなら \(c\,e^{-i\varphi}=|c|\)
となり \(y\equiv0\) で読出しは退化する）。\(\varphi\)
は親平面から独立に読まれる集団時計位相であり、\(c\) は \(Z_\perp\)
の座標である------二つは別の読出しであり、\(c\,e^{-i\varphi}\)
は\textbf{時計に対する相対複素座標}を与える。\(x\) と \(y\)
の直交性は、複素数の実部と虚部の直交性------すなわち\textbf{読出しの構成}------であって、実体側の二本の直交軸に対応しない（実体は1複素自由度である）。第3軸は回転軸方向の復調射影

\[z(n) = \mathrm{Re}\,\big(\langle \hat a, Z_\perp\rangle\, e^{-i\varphi}\big)\]

である。等方に見える直交3軸は、こうして「複素1つの位相分解＋従属軸」として製造される。

\textbf{実測}（図2）：時計に同期した成分は一定半径 \(|r|=0.983\pm0.005\)
のまま角度のみ緩慢に歳差する（0.12\%/時計）------振幅不変の準静止。静止成分への変位（キック）は8000ステップ無減衰で保持される。

\begin{figure}
\centering
\pandocbounded{\includegraphics[keepaspectratio,alt={図2 quadrature読出し}]{figs/fig_s2_quadrature_v1.png}}
\caption{図2 quadrature読出し}
\end{figure}

\textbf{図2}：(a) 時計同期成分の準静止。(b) 変位の保持。

先行論文が外から与えた3次元レジスタ格子は、この内在読出しの外部代用だったことが分かる。

\subsection{7. 固有時間 τ
の読出しの導出}\label{ux56faux6709ux6642ux9593-ux3c4-ux306eux8aadux51faux3057ux306eux5c0eux51fa}

\textbf{導出。} 錐の右辺 \(R'^2 = t^2+R^2+Q^2\)
に対応する時計として、局所系が読めるのは合計 \(R'\)
の刻みである（内訳は§10）。凝縮体においてそれは\textbf{親平面の回転位相
\(\varphi(n)\)}------単一の集団時計------として現れる。固有時間の読出しは

\[\tau(n) = \varphi(n)\,/\,\omega_{\text{clock}}\]

である。

\textbf{実測1：時計は1本（図3a）。}
高分解能測定（窓10倍・ゼロパディング）で、全占有平面の回転周波数は単一値
\(\omega \approx \pi/72\)/step（力学の基本角と0.04\%一致）に収束する（比
\(1\pm1\%\)）。周波数の梯子は存在しない（低分解能で見えた「梯子」は計器アーティファクト------§13）。

\textbf{実測2：時計は引き込みで集団強制される（図3b）。}
凝縮体の自然揺らぎ成分は位相スリップを含む非同期の位相を示すが、同じ方向への\textbf{コヒーレントな}キックは、固有の離調を保持できず凝縮体の時計にロックする（ドリフト
\(-1.4\% \to -0.05\%\)）。固有時間は個々の成分の属性ではなく、\textbf{凝縮体が引き込みによって集団強制する秩序}である。

\begin{figure}
\centering
\pandocbounded{\includegraphics[keepaspectratio,alt={図3 単一時計と引き込み}]{figs/fig_s3_clock_entrainment_v1.png}}
\caption{図3 単一時計と引き込み}
\end{figure}

\textbf{図3}：(a) 全占有平面が単一時計で剛体回転（4つのN）。(b)
引き込み------摂動は凝縮体の固有時計に同期させられる。

\textbf{帰結：静止と運動の分類。}
静止＝時計ロック（変位は保持される、図2b）。時計を持たない摂動はすべて引き込まれて静止する。ゆえに\textbf{運動は、引き込みに抗して自分の位相を刻み続ける内容------自前の時計を持つもの＝質量------の特権}である。質量側の実証は先行論文の物質パケット並進
{[}1{]}
が既に与えている------本稿の読出し語彙で全体を可視化したのが図8である：xyz読出しの3次元軌道は等間隔の点列（等速並進）を描き、x--τ・y--τ・z--τ
の3投影面すべてで実測点が予言直線に乗る（τ=0〜900、10点全て偏差0.000、PR一定3.39------900ステップの間、形を一切変えずに直進する）。質量→固有時計→時間→運動の系列が閉じる。

\begin{figure}
\centering
\pandocbounded{\includegraphics[keepaspectratio,alt={図8 4D並進}]{figs/fig_s8_translation4d_v1.png}}
\caption{図8 4D並進}
\end{figure}

\textbf{図8}：xyz＋τ（R′固有時間）の4次元読出し。(a)
3次元軌道------等間隔の点＝等速並進。(b)(c)(d) x--τ, y--τ, z--τ
面への投影------直線＝等速。全10点で偏差0.000。

\subsection{8. 質量的な量 m
の読出しの再導出}\label{ux8ceaux91cfux7684ux306aux91cf-m-ux306eux8aadux51faux3057ux306eux518dux5c0eux51fa}

質量は R
成分の大きさであり、波形の\textbf{構造}（スナップショットに写る量）から読める。3系統を導出する。

\subsubsection{8.1
読出し(i)：非コヒーレンス（全段導出）}\label{ux8aadux51faux3057iux975eux30b3ux30d2ux30fcux30ecux30f3ux30b9ux5168ux6bb5ux5c0eux51fa}

二成分の複素読出し
\(u=(c_1, c_2)\)（例：親平面座標と平面外主方向座標）の窓平均 Gram 行列

\[\Gamma = \langle u\,u^\dagger\rangle = \begin{pmatrix} \langle|c_1|^2\rangle & \langle c_1\bar c_2\rangle \\ \langle c_2\bar c_1\rangle & \langle|c_2|^2\rangle \end{pmatrix}\]

は \(2\times2\) エルミートである。Pauli 基底で分解すると
\(\Gamma = T\,\mathbb{1} + X\sigma_x + Y\sigma_y + Z\sigma_z\)、行列式は恒等的に

\[\det\Gamma = T^2 - X^2 - Y^2 - Z^2.\]

\(\Gamma\) は半正定値だから
\(\det\Gamma \ge 0\)------\textbf{光錐型の束縛が自動で成立する}。等号
\(\det\Gamma=0\) は \(\Gamma\) が rank 1、すなわち窓内で \(u\)
が単一の複素振幅に比例する（\(c_2 = \text{const}\cdot c_1\)）------完全コヒーレント＝\textbf{光的}------のとき、かつそのときに限る。そこで

\[m^2 \equiv \det\Gamma\]

を質量の読出しとする：質量とは読出し対の\textbf{非コヒーレンス}である。規格化を区別しておく：\(\det\Gamma\)
は振幅の4次量なので、絶対的な質量候補は
\(\sqrt{\det\Gamma}\)、無次元の\textbf{質量度}は
\(\det\Gamma/T^2 \in [0,1]\)
である。本稿の実測（§9.5を含む）が測るのは後者であり、「最大質量」とは正確には最大非コヒーレンス度である------絶対エネルギースケール（R軸の尺度）の確定は将来課題として残る。

\textbf{数学の出所の明示。} この数学は新しくない：\(2\times2\)
コヒーレンス行列の Pauli 分解係数は偏光の Stokes
パラメータであり、\(\det J \ge 0\)（等号＝完全偏光）は偏光理論の古典的恒等式である
{[}11,12{]}。本稿の寄与は物理的\textbf{同定}------非偏光度を質量²と読み、光的＝完全偏光（rank
1）とする辞書------と、次の模型内実測である。

\textbf{機械検証}（図4、判定は実行前固定）：rank-1合成状態で
\(m^2/T^2 = 1.4\times10^{-16}\)（機械零）、独立2成分の非コヒーレント合成で
\(1.000\)。模型実測では、凝縮前の親固有モード（光的候補）が
\(1.2\times10^{-7}\)、準安定凝縮体が
\(3.1\times10^{-3}\)------\textbf{凝縮体は光的固有モードの
\(2.5\times10^4\)
倍の質量を持つ}。§7の引き込み（凝縮体は自前の時計を強制する）と整合する。

\begin{figure}
\centering
\pandocbounded{\includegraphics[keepaspectratio,alt={図4 質量=非コヒーレンス}]{figs/fig_s6_mass_check_v1.png}}
\caption{図4 質量=非コヒーレンス}
\end{figure}

\textbf{図4}：質量読出し＝Gram行列の非コヒーレンス。4状態の機械検証。

\subsubsection{8.2
読出し(ii)：閉鎖完走コスト}\label{ux8aadux51faux3057iiux9589ux9396ux5b8cux8d70ux30b3ux30b9ux30c8}

生成を「閉鎖の完走」と読む枠（第9論文
{[}7{]}）では、質量は完走に要するコストとして読める（操作的定義と
\(f(N)=2\ln N/N\) 型の計数は {[}7{]}
が与える。同定の実証は同論文の実験系列）。本稿では出所表示に留め、主読出しは
(i)(iii) とする。

\subsubsection{8.3
読出し(iii)：自前時計の強度}\label{ux8aadux51faux3057iiiux81eaux524dux6642ux8a08ux306eux5f37ux5ea6}

§7の帰結から、質量は\textbf{引き込みに抗する度合い}------自前の位相を刻み続ける強度------として力学的に読める。光的な線形摂動は完全に引き込まれ（質量ゼロ）、物質内容は自分の時計を保って運動する。読出し(i)との整合：完全コヒーレント（rank
1）⟺ 引き込まれる ⟺ 質量ゼロ、の対応が §7・図4
の実測で同時に成立している。

\subsection{9. 電荷的な量 q
の読出しの再導出}\label{ux96fbux8377ux7684ux306aux91cf-q-ux306eux8aadux51faux3057ux306eux518dux5c0eux51fa}

電荷は Q
成分の大きさ＋相対符号であり、同じく波形の構造から読める。3系統の操作的定義を明記する。

\textbf{読出し(i)：透過率。} 二波の衝突読出し
\(\theta = \operatorname{atan2}(\sqrt{P_f}, \sqrt{P_b})\)（合成スペクトルのフェルミオン/ボゾン関係量比）から反射率
\(R=\sin^2\theta\) が読め、その補である透過率 \(1-R\)
が素電荷の読出しである（素電荷との定量同定 \(1-R=\sqrt{4\pi\alpha}\)
の実証は三部作A {[}6{]}）。

\textbf{読出し(ii)：相対位相（電荷文法）。}
電荷の\textbf{符号}は二波の相対位相として読める（同符号＝一方の位相関係、異符号＝反転。中性化の二機構を含む文法の実証は
{[}6{]}）。§11の原理（物理量は相対符号のみ）と整合する------電荷符号も単体属性ではなく相対量である。

\textbf{読出し(iii)：巻き数。}
内部軸（\(\eta\)）方向のフーリエ変換の支配巻き数 \(m\)
が電荷型（KK型）の運動量である。対生成が巻き数の和則
\(m^* = 2m_B - m_s\) を厳密に守ることは二体・多体のcensus {[}2,1{]}
で実測済み（予言外の毛は機械零）。

三系統はいずれも波形のスナップショットから読める構造量であり、閉塞（二乗しか見ない）と整合して、読めるのは大きさと相対符号------符号つき座標そのものではない。

\subsection{9.5
読出しの分布------海の中の質量・電荷・スピン・分類（実測の補遺）}\label{ux8aadux51faux3057ux306eux5206ux5e03ux6d77ux306eux4e2dux306eux8ceaux91cfux96fbux8377ux30b9ux30d4ux30f3ux5206ux985eux5b9fux6e2cux306eux88dcux907a}

§8--§9の読出しを、三つの状態------対称な海の成長期（j=400--500）・熱平衡期（j=3000--3200）・帯電した構造（census型：単一巻きポンプ＋単一巻き種）------のモード集合
\((k, m)\)（生bin×巻き数）に適用し、分布を実測した。

\textbf{質量の分布（図9）}：海のモードは、成長期・熱平衡期のどちらでも\textbf{最大非コヒーレンス（質量²/T²≈1）に集中}する------海は「重い」。軽い（光的）状態は海には現れず、コヒーレントな構造（固有モード
\(1.2\times10^{-7}\)、図4）に限られる。質量スペクトルの構造は海ではなくコヒーレント構造が担う。

\textbf{電荷の分布（図10）}：対称な海では電荷（巻き数）は多数の値にほぼ平坦に分布し、\textbf{±1
に安定しない------これは既知の未解決課題であり、本実測はその分布の直接測定である}（§15-6）。帯電した構造では複数の電荷値が共存する：種
\(m=1\) が支配し、和則相棒 \(m=3\)
が機械ゼロの床から明瞭に立つ（対数表示）------電荷は量子化された値の\textbf{集合}を取るが、±1への収斂機構は未発見である。

\textbf{ボゾン的な波とフェルミオン的な波の分布と、スピン的な量の分布（図11）}：クラス別パワースペクトルは、成長期の構造（ポンプ線・フェルミオン的な山）が熱化で平坦化していく過程を示す（F/B総和
21.9/42.2 → 28.9/35.1------等分配への接近）。スピン的な量（Gram の Bloch
ベクトル------§8.1の副産物としての\textbf{探索的}読出し）は、海ではほぼ無偏極（\(|S|\approx0.03\)）、帯電構造のコヒーレント種は約5倍偏極（\(|S|\approx0.13\)）だが完全偏極には遠い。スピンの本格的同定は将来課題である。

\textbf{粒子的・反粒子的・灰色状態の分布（図12）}：フェルミオン的な波の符号均衡
\(b=(P_+-P_-)/(P_++P_-)\) の分布は、対称な海では \(b=0\)
近傍（\textbf{灰色}）に集中し（成長期・熱平衡期とも粒子率・反粒子率ゼロ）、帯電した種を持つ構造では
\(b=+1\)（\textbf{粒子的}）に集中する。粒子/反粒子/灰色は状態の三分類として読出し可能であり、対称な初期条件からは灰色の海が、帯電した種からは粒子的な内容が育つ。

\begin{figure}
\centering
\pandocbounded{\includegraphics[keepaspectratio,alt={図9 質量分布}]{figs/fig_s9_mass_distribution_v1.png}}
\caption{図9 質量分布}
\end{figure}

\textbf{図9}：質量として読み出せる量の分布。海は最大非コヒーレンスに集中する。

\begin{figure}
\centering
\pandocbounded{\includegraphics[keepaspectratio,alt={図10 電荷分布}]{figs/fig_s10_charge_distribution_v1.png}}
\caption{図10 電荷分布}
\end{figure}

\textbf{図10}：電荷として読み出せる量の分布。±1への安定化は未解決（オープン課題）。

\begin{figure}
\centering
\pandocbounded{\includegraphics[keepaspectratio,alt={図11 クラスとスピン}]{figs/fig_s11_class_spin_v1.png}}
\caption{図11 クラスとスピン}
\end{figure}

\textbf{図11}：ボゾン的・フェルミオン的な波の分布と、スピン的な量の分布。

\begin{figure}
\centering
\pandocbounded{\includegraphics[keepaspectratio,alt={図12 粒子・反粒子・灰色}]{figs/fig_s12_particle_gray_v1.png}}
\caption{図12 粒子・反粒子・灰色}
\end{figure}

\textbf{図12}：フェルミオン的な波の粒子的・反粒子的・灰色状態の分布。

\subsection{10.
中心命題------観測の全内容と、唯一読めないもの}\label{ux4e2dux5fc3ux547dux984cux89b3ux6e2cux306eux5168ux5185ux5bb9ux3068ux552fux4e00ux8aadux3081ux306aux3044ux3082ux306e}

閉塞の錐

\[x^2 + y^2 + z^2 \;=\; t^2 + R^2 + Q^2 \;=\; R'^2\]

の上に、§6--§9 の読出しを整理する（図5）。

{\def\LTcaptype{none} % do not increment counter
\begin{longtable}[]{@{}
  >{\raggedright\arraybackslash}p{(\linewidth - 4\tabcolsep) * \real{0.3333}}
  >{\raggedright\arraybackslash}p{(\linewidth - 4\tabcolsep) * \real{0.3333}}
  >{\raggedright\arraybackslash}p{(\linewidth - 4\tabcolsep) * \real{0.3333}}@{}}
\toprule\noalign{}
\begin{minipage}[b]{\linewidth}\raggedright
量
\end{minipage} & \begin{minipage}[b]{\linewidth}\raggedright
読出し
\end{minipage} & \begin{minipage}[b]{\linewidth}\raggedright
手段（本文の導出）
\end{minipage} \\
\midrule\noalign{}
\endhead
\bottomrule\noalign{}
\endlastfoot
\(x, y, z\) & ○ 直接 & quadrature（§6） \\
固有時間 τ = R′時計 & ○ 直接 &
単一時計の刻み・引き込みで集団強制（§7） \\
質量 m（R成分） & ○ 分離可 &
非コヒーレンス／完走コスト／自前時計強度（§8） \\
電荷 q（Q成分） & ○ 分離可 & 透過率／相対位相／巻き数（§9） \\
\textbf{座標時間 t} & \textbf{× 唯一不可} & 残差 \(t^2=R'^2-m^2-q^2\)
または広域選択のみ \\
\end{longtable}
}

\begin{figure}
\centering
\pandocbounded{\includegraphics[keepaspectratio,alt={図5 閉塞の錐の読出し構造}]{figs/fig_s7_cone_readout_v1.png}}
\caption{図5 閉塞の錐の読出し構造}
\end{figure}

\textbf{図5}：何が読めて、何が読めないか。

\textbf{唯一読めないのは座標時間 t である。}
本シリーズの全実験系を見渡しても、t
の直接読出し器は存在しない（実験の「時間」は常にステップカウンタ＝規約か、R′時計である）。t
が現れる経路は二つしかない：(a) 算術残差 \(t^2 = R'^2 - m^2 - q^2\)、(b)
大きな系の空間振動群（多数の局所時計）を突き合わせ、「広域に共有される時間がある」と想定して\textbf{選択}する経路------実際、日常の
t の分離はこの経路であり、R や Q
の分離も空間内のバネ振動・電磁振動の観測からの帳簿分離である。

\textbf{構造的理由の候補}：R と Q
の内容は波形の\textbf{構造}（非コヒーレンス・巻き数------スナップショットに写る）だが、t
の内容は\textbf{流れ}そのものであり、スナップショットに痕跡を残さない。流れの中から流れは読めない。

\textbf{帰結1：固有時間／座標時間の区別の導出。}
一般相対論では「固有時間は時計が測る物理量・座標時間はゲージ（規約）」の区別が微分同相不変性という\textbf{原理}から従う。本稿ではこの区別が、閉塞の錐の読出し構造から\textbf{導出}される------原理の向きが逆である。

\textbf{帰結2：なぜ 3＋1 か。} 観測の構造は厳密に
3＋1（実3軸＋時計1本）＋物質属性（m,
q）である。読出し可能な内容がちょうどそれだから、というのが本模型の答えの候補である。

\subsection{11.
向きの規約------時間の矢は存在しない}\label{ux5411ux304dux306eux898fux7d04ux6642ux9593ux306eux77e2ux306fux5b58ux5728ux3057ux306aux3044}

時間の矢は本模型に存在しない。虚数符号は「この軸の向きは読めない」という記帳であり（\((it)^2=(-it)^2\)------閉塞は二乗しか見ない）、写像を前向きに回すことは物理ではなく規約である。実測もこれと整合する：初期状態を複素共役に置き換えても（状態の共役は時間反転ではない------§13）、規約不変量である\textbf{相対符号}------時計の回転と回転軸の向きの相対掌性------は不変（−）に留まった。物理量は相対符号のみであり、絶対の向き（時間の矢・空間の掌性）は選択問題である。なお本節の主張は「基本閉塞と可逆写像に絶対的な時間向きがない」ことであり、熱化・生成・自己触媒の履歴から\textbf{統計的な有効時間の矢}が創発するか否かは別の問いである（§15-9）。

\subsection{12.
適用域の境界------次元の鋭さと少数体領域}\label{ux9069ux7528ux57dfux306eux5883ux754cux6b21ux5143ux306eux92edux3055ux3068ux5c11ux6570ux4f53ux9818ux57df}

第3次元の従属性（回転軸＝外積）の一致度は完全ではなく、Nとともに向上する。曖昧さ
\(1-|\hat n\cdot\hat a|\) は関係波の本数 \(M\)
の逆数程度でスケールし（0.0126 → 0.0082 → 0.0035、N=5,6,8）、\textbf{N=4
では曖昧さ 0.9986（一致度
\(|\hat n\cdot\hat a|=0.0014\)）------第3次元がそもそも結晶化しない}（図6）。段階を明示すると：N=4＝非結晶化、N=5＝結晶化するが曖昧さ最大（少数体領域の縁）、大N＝鋭い空間相、という連続的な階層である。

\begin{figure}
\centering
\pandocbounded{\includegraphics[keepaspectratio,alt={図6 次元の鋭さ}]{figs/fig_s4_sharpness_v1.png}}
\caption{図6 次元の鋭さ}
\end{figure}

\textbf{図6}：次元の曖昧さ \textasciitilde1/M。N=4
では第3次元が結晶化しない。

これは欠陥ではなく適用域の境界である。N=4,5
は閉じた関係系としては\textbf{少数体複合粒子------バリオン------の領域}であり、極端に量子化・離散化され、エイリアスが多く、次元も時間も存在すら曖昧なのはその領域の物理そのものである（クォークの分数電荷＝低位数の有理住所という本シリーズの別系列とも整合する）。本稿の対象は空間が存在する大N側であり、少数体領域は別の研究主題として明示的に切り出す。

\subsection{13.
反証と修正の記録}\label{ux53cdux8a3cux3068ux4feeux6b63ux306eux8a18ux9332}

{\def\LTcaptype{none} % do not increment counter
\begin{longtable}[]{@{}
  >{\raggedright\arraybackslash}p{(\linewidth - 6\tabcolsep) * \real{0.1429}}
  >{\raggedright\arraybackslash}p{(\linewidth - 6\tabcolsep) * \real{0.2857}}
  >{\raggedright\arraybackslash}p{(\linewidth - 6\tabcolsep) * \real{0.2857}}
  >{\raggedright\arraybackslash}p{(\linewidth - 6\tabcolsep) * \real{0.2857}}@{}}
\toprule\noalign{}
\begin{minipage}[b]{\linewidth}\raggedright
\#
\end{minipage} & \begin{minipage}[b]{\linewidth}\raggedright
誤り
\end{minipage} & \begin{minipage}[b]{\linewidth}\raggedright
発見
\end{minipage} & \begin{minipage}[b]{\linewidth}\raggedright
修正
\end{minipage} \\
\midrule\noalign{}
\endhead
\bottomrule\noalign{}
\endlastfoot
1 & 位相勾配計：復調位相の傾きで周波数を測定 &
振幅のゼロ交差で位相スリップ・汚染 &
FFTピーク計へ。基線の「離調」数値はスリップ汚染を含むため、引き込みの結論は「スリップする非同期
→ スリップなしのロック」として表現（図3b） \\
2 & \textbf{FFTビン量子化を「0.1\%精度の等間隔コム」と誤読} &
「間隔」が全NでFFTビン幅 \(2\pi/n_s\) と厳密一致 &
窓10倍＋ゼロパディングで再測定→コム消滅・単一時計（図7）。\textbf{分解能の限界は構造をでっち上げる------本稿の主題が計器側に現れた実例} \\
3 & 窓過渡：早期窓の周波数ドリフトを実構造と誤認 &
高分解能・後期窓で消滅 & 整定後の窓に統一 \\
4 & 状態の複素共役を時間反転と混同（キラリティ検定の設計誤り） &
共役状態も同符号へ収束 &
「物理量は相対符号のみ」への再解釈（§11）。共役検定は規約不変量の確認として成立 \\
5 & 「運動＝離調」の素朴設計 & キックは離調を保持せず時計にロック &
\textbf{引き込みの発見}（§7）------反証が中心結果に転じた \\
\end{longtable}
}

\begin{figure}
\centering
\pandocbounded{\includegraphics[keepaspectratio,alt={図7 偽コム}]{figs/fig_s5_false_comb_v1.png}}
\caption{図7 偽コム}
\end{figure}

\textbf{図7}：計器教訓\#2。低分解能の「等間隔コム」（左）はFFTビン幅そのものであり、高分解能（右）で消滅する。

5件のうち3件は「測るべき量の性質を導出せずに計器を作った」誤りである。規律として明文化する：\textbf{計器は、測定量の不変性（または予言値）を導出してから設計する。}

\subsection{14.
先行研究------独自導出の後に見出した対応}\label{ux5148ux884cux7814ux7a76ux72ecux81eaux5c0eux51faux306eux5f8cux306bux898bux51faux3057ux305fux5bfeux5fdc}

本稿の読出し構成は閉塞の錐と凝縮体の観測量からの独自導出であり、以下に依拠して設計したものではない。導出後の再調査により次の対応を見出した。対応は傍証であり、証明ではない。

\begin{enumerate}
\def\labelenumi{\arabic{enumi}.}
\tightlist
\item
  \textbf{関係論的時間------Page--Wootters機構}
  {[}8{]}。大域的に定常な状態の中で、時計部分系との相関から時間発展が現れる（evolution
  without
  evolution）。相違：あちらは量子的制約ともつれからの条件付け。本稿は時計そのものが力学的に自己生成され（凝縮＋引き込み）、座標時間の非可読性が読出し構造から導出される。
\item
  \textbf{座標時間はゲージ------部分観測量}
  {[}9{]}。一般相対論で座標時間が観測量でないことは正統な既知事実であり、Rovelli
  の partial observables
  が定式化した。相違：GRでは微分同相不変性が原理として課され、その帰結として座標時間が非物理になる。本稿では読出し構造からこの区別が導出される------\textbf{原理の向きが逆}。
\item
  \textbf{quadrature読出し------ホモダイン検波}
  {[}10{]}。局部発振器（LO）の位相を基準に直交位相成分を読む量子光学の標準技術は、本稿の
  \(x, y\)
  読出しと同一の数理形式である。相違：光学ではLOは実験者が供給する外部装置。本稿では位相基準が系自身の凝縮体（引き込みで自己選択された固有時計）------\textbf{宇宙が自分のLOを供給する}。
\item
  \textbf{質量読出しの数学------偏光のコヒーレンス行列}
  {[}11,12{]}。\(2\times2\) コヒーレンス行列の Pauli 分解＝Stokes
  パラメータ、\(\det J\ge0\)（等号＝完全偏光）は偏光理論の古典。本稿の寄与は物理的同定（非偏光度＝質量²・光的＝完全偏光）と模型内実測に限られる。
\item
  \textbf{次元のスケール依存------CDTのスペクトル次元}
  {[}13{]}。短距離で時空の実効次元が減る（running dimensional
  reduction）。本稿の「曖昧さ\textasciitilde1/M・少数体領域で非結晶化」と同方向。相違：あちらは拡散次元、こちらは回転軸と外積の一致度という構成的次元計。
\item
  \textbf{引き込み------同期理論}
  {[}14{]}。結合振動子の集団リズムへの引き込みは Huygens
  以来の古典。相違：本稿は引き込みを\textbf{固有時間の存立機構}として読み、「静止は無質量摂動の運命・運動は質量の特権」という運動学の分類へ接続する。
\item
  \textbf{物理系としての基準系------量子基準系}
  {[}15{]}。方向や時刻は物理系に相対的にのみ定義される（unspeakable
  information）。相違：あちらは情報理論的資源論で基準系は与えられる。本稿では基準系（時計と軸）が力学から自己生成されることを実測する。
\end{enumerate}

\subsection{15.
展望------未解決課題（進行指針、本稿は主張しない）}\label{ux5c55ux671bux672aux89e3ux6c7aux8ab2ux984cux9032ux884cux6307ux91ddux672cux7a3fux306fux4e3bux5f35ux3057ux306aux3044}

\begin{enumerate}
\def\labelenumi{\arabic{enumi}.}
\tightlist
\item
  \textbf{質量つき内容の運動の内在座標版}：引き込みに抗して動く物質パケットを、外置きレジスタでなく§6の内在座標で追う（「運動＝質量の特権」の直接完結）。
\item
  \textbf{t選択の力学}：大きな系に複数の局所凝縮体を置き、広域時間 t
  がどう選択されるか（時計たちの突き合わせ＝同期ネットワーク）を実測する。
\item
  \textbf{枠と計量の完全版}：弦長公式 \(R''=2R\sin(\Delta\theta/2)\)
  による計量の構成、曲率、物質の逆作用。
\item
  \textbf{少数体（バリオン）領域の専用調査}：次元・時間・存在が曖昧な領域の物理------別テーマとして。
\item
  \textbf{±1\%残差の帰属}：分解能か微細構造か（優先度低として記録）。
\item
  \textbf{電荷の±1安定化}：電荷（巻き数）は複数の量子化値に分布し、±1への収斂機構は未発見（図10）。素電荷の一意性がどこから来るかは、本シリーズの主要オープン課題である。あわせて三読出し（透過率・相対位相・巻き数）を同一状態集合で同時測定し、単調関係・符号対応・保存・対生成時の共役対応を一つの表に統一することも残る。
\item
  \textbf{質量度と引き込み耐性の連続掃引}：非コヒーレンス度
  \(\mu=\det\Gamma/T^2\)
  を0から1まで掃引し、同期時間・残留離調・自前時計・速度を測る。\(\mu\)
  と引き込み耐性が一価関数で結ばれれば、「質量＝非コヒーレンス」は同定から力学的質量へ進む。総和
  \(\sum \det\Gamma\) が保存量・単調量のいずれに当たるかも同時に調べる。
\item
  \textbf{読出し代数の完全性}：許容される不変読出し（瞬時スナップショット・有限窓相関、共通位相・置換・規約不変）を分類し、その不変量環が
  \(\{x,y,z,\tau,m,q\}\) で生成され独立な \(t\)
  を含まないことを示す------「唯一不可読」の no-go 定理化。
\item
  \textbf{統計的な時間の矢}：基本力学に絶対向きはない（§11）が、熱化・自己触媒の履歴からエントロピー的な有効矢が創発するかは未検討の別課題。
\end{enumerate}

\subsection{16.
本稿が主張しないこと・反証条件}\label{ux672cux7a3fux304cux4e3bux5f35ux3057ux306aux3044ux3053ux3068ux53cdux8a3cux6761ux4ef6}

本稿は、現実の時空がこの模型であるとは主張しない。反証条件：(a)
許容される読出し（波形スナップショットまたは有限窓相関の、共通位相・ノード置換・規約に不変な関数）から、\(\tau, m, q\)
と独立で座標変換の選択に依らない\textbf{時間読出し量 t
が構成されれば}、「唯一不可読」命題は棄却される------公開コード上で探索・反証可能な形である。(b)
時計を持たない摂動が引き込みを受けずに運動する例。(c)
読出しに依らず実体側に直交3軸が実在する証拠（同率占有の独立な複素固有対の共存等）。(d)
曖昧さの 1/M 則の破れ。(e)
質量読出し(i)(iii)の乖離（非コヒーレンスと引き込み耐性が独立に振る舞う例）。

\textbf{結び}：凝縮体は、位置の軸と時計を自分で供給した。軸は複素1つの位相分解として、時計は引き込みが強制する単一の刻みとして。そして読めない量がただ一つ残った------座標時間。固有時間は物理であり、座標時間は選択である------この区別ごと、閉塞の錐の読出し構造から出てきた。

\subsection{17. 再現性}\label{ux518dux73feux6027}

シリーズ規約（本稿単体で検証が閉じる・原本は read-only
import・判定は実行前固定）に従う。実験スクリプト8本（4D並進読出し
\texttt{run\_translation\_4d\_readout\_v1.py}・分布読出し3系
\texttt{run\_distribution\_readouts\_v1.py}
を含む）・図生成2本・結果JSON 10点＋npz 1点の対応表は同フォルダ
\texttt{実験一覧\_v1.md}
を正本とする。全スクリプトは決定論であり、保存JSONとの再現一致を機械確認済み（v3b/v3c
で実施）。系列を保存しない旧版実験（v4）は同一シード再走行＋要約値対照検定（6桁一致）の上で図に使用した。計器の失敗2件（位相勾配計・ビン量子化）はスクリプトと結果JSONごと保存されている。

\subsection{参考文献}\label{ux53c2ux8003ux6587ux732e}

\textbf{シリーズ内（出所表示。本文の論理はいずれにも依存しない------全読出しは本文で導出済み）}

{[}1{]}
木原範昭『インフレーションの終わり方が、物質生成の始まり方である』(2026),
Zenodo. Concept DOI:
10.5281/zenodo.21809814（先行論文：外置きの位置軸と時計・物質パケット並進・毛のcensus多体版）

{[}2{]}
木原範昭『フェルミオン的構造の生成は誘導・自己触媒・対相関で起こる』(2026),
Zenodo. Concept DOI:
10.5281/zenodo.21808091（2＋1構造の先行形＝閉形式回転・読出し側分類の原理・census二体版）

{[}3{]} 木原範昭『第7論文：N体関係波閉鎖系の時間構造』(2026), Zenodo.
Concept DOI: 10.5281/zenodo.21543070（三方向準安定構造の発見）

{[}4{]} 木原範昭『第8論文：二段階seed除去による因果分離』(2026), Zenodo.
Concept DOI: 10.5281/zenodo.21614402（力学コードの直接原本）

{[}5{]} 木原範昭『無名等振幅複合波モデル 基本公理系』(2026), Zenodo.
Concept DOI: 10.5281/zenodo.21315735（零二乗和閉塞------閉塞の錐の公理）

{[}6{]} 木原範昭『相互作用の二文法分解』(Anonymous Two-Channel Closed
Wave System Trilogy, Part I) (2026), Zenodo. Concept DOI:
10.5281/zenodo.21763995（透過率＝素電荷の同定・電荷文法の実証）；『AB二体閉鎖位相系における未来位相位置加速度写像と調和閉鎖による逆二乗則』(2026),
Zenodo. DOI: 10.5281/zenodo.21441081（先行論文の外置き分散時計の出所）

{[}7{]} 木原範昭『フェルミオンの生成構造』(2026), Zenodo. Concept DOI:
10.5281/zenodo.21766706（閉鎖完走コストの操作的定義）

\textbf{外部}

{[}8{]} D. N. Page, W. K. Wootters, ``Evolution without evolution:
Dynamics described by stationary observables,'' Phys. Rev.~D
\textbf{27}, 2885 (1983)

{[}9{]} C. Rovelli, ``Partial observables,'' Phys. Rev.~D \textbf{65},
124013 (2002)

{[}10{]} U. Leonhardt, \emph{Measuring the Quantum State of Light}
(Cambridge Univ. Press, 1997)；H. P. Yuen, V. W. S. Chan, ``Noise in
homodyne and heterodyne detection,'' Opt. Lett. \textbf{8}, 177 (1983)

{[}11{]} G. G. Stokes, ``On the composition and resolution of streams of
polarized light from different sources,'' Trans. Cambridge Philos. Soc.
\textbf{9}, 399 (1852)

{[}12{]} E. Wolf, ``Coherence properties of partially polarized
electromagnetic radiation,'' Nuovo Cimento \textbf{13}, 1165 (1959)；M.
Born, E. Wolf, \emph{Principles of Optics}, 7th ed.~(Cambridge Univ.
Press, 1999)

{[}13{]} J. Ambjørn, J. Jurkiewicz, R. Loll, ``Spectral dimension of the
universe,'' Phys. Rev.~Lett. \textbf{95}, 171301 (2005)

{[}14{]} A. Pikovsky, M. Rosenblum, J. Kurths, \emph{Synchronization: A
Universal Concept in Nonlinear Sciences} (Cambridge Univ. Press, 2001)

{[}15{]} S. D. Bartlett, T. Rudolph, R. W. Spekkens, ``Reference frames,
superselection rules, and quantum information,'' Rev.~Mod. Phys.
\textbf{79}, 555 (2007)

\end{document}
